\section{Conclusion}\label{sec:conclusion}
In this paper, we present \ours, a critic-free RLHF framework designed to enhance training stability by addressing fundamental flaws in existing methods. We identified that prior critic-free algorithms, such as GRPO, rely on \textit{prompt-level (local)} advantage normalization, which we prove to be theoretically biased (Appendix \Cref{appendix:proof_GRPO}) and demonstrate to be practically unstable.

Our solution, \textbf{Global Advantage Normalization}, normalizes advantages across the entire batch, providing a stable and \textit{effectively unbiased} estimator (whose bias vanishes as batch size increases). We proposed two variants tailored for different needs:
\begin{enumerate}
    % --- MODIFIED: Conclusion updated to match the k>=1 logic ---
    \item \ours: For general-domain RLHF, we show the single-sample ($k=1$) approach is highly efficient and achieves superior OOD generalization. We also show that its application with $k>1$ is robust for reasoning tasks (Section 4.2).
    % --- END OF MODIFICATION ---
    \item \varn: For complex reasoning and agent tasks, this variant combines group-mean reshaping with stable global normalization and a theoretically sound $k_2$ estimator for KL.
\end{enumerate}

Our empirical results, supported by independent validation from multiple third-party systems, confirmed the effectiveness of this specialized approach. \ours demonstrated state-of-the-art generalization in general RLHF. \varn showed dramatic improvements in stability, preventing overfitting in low-data regimes and outperforming both GRPO and full-critic PPO in complex, long-horizon tool-use tasks. This work demonstrates that a theoretically sound, stable approach without a critic model can be more efficient and effective than traditional PPO.
